\documentclass{article}
\usepackage{apacite}
\usepackage{hyperref}
\usepackage{multicol}
\usepackage{natbib}
\usepackage{fiona}
%\usepackage{elijah}
\usepackage[utf8]{inputenc}
\usepackage{array}
\setlength{\parindent}{0pt}
\usepackage{hyperref}

\usetikzlibrary{calc,patterns,angles,quotes}
\usepackage{pgfplots}
\pgfplotsset{width=15cm,height=10cm,compat=1.9}

\usepackage[english]{babel}
\usepackage{fancyhdr}
\usepackage{lastpage}
\pagestyle{fancy}
\fancyhf{}


\pagestyle{empty}% for cropping
\usepackage{tikz}
\usetikzlibrary{calc}
\usepackage{zref-savepos}
\usepackage{tabu}

\newcounter{NoTableEntry}
\renewcommand*{\theNoTableEntry}{NTE-\the\value{NoTableEntry}}

\newcommand*{\strike}[2]{%
  \multicolumn{1}{#1}{%
    \stepcounter{NoTableEntry}%
    \vadjust pre{\zsavepos{\theNoTableEntry t}}% top
    \vadjust{\zsavepos{\theNoTableEntry b}}% bottom
    \zsavepos{\theNoTableEntry l}% left
    \hspace{0pt plus 1filll}%
    #2% content
    \hspace{0pt plus 1filll}%
    \zsavepos{\theNoTableEntry r}% right
    \tikz[overlay]{%
      \draw
        let
          \n{llx}={\zposx{\theNoTableEntry l}sp-\zposx{\theNoTableEntry r}sp-\tabcolsep},
          \n{urx}={\tabcolsep},
          \n{lly}={\zposy{\theNoTableEntry b}sp-\zposy{\theNoTableEntry r}sp},
          \n{ury}={\zposy{\theNoTableEntry t}sp-\zposy{\theNoTableEntry r}sp}
        in
        (\n{llx}, \n{lly}) -- (\n{urx}, \n{ury})
      ;
    }% 
  }%
}

\usepackage{tikz}
\usetikzlibrary{calc}

\tikzset{
    right angle quadrant/.code={
        \pgfmathsetmacro\quadranta{{1,1,-1,-1}[#1-1]}     % Arrays for selecting quadrant
        \pgfmathsetmacro\quadrantb{{1,-1,-1,1}[#1-1]}},
    right angle quadrant=1, % Make sure it is set, even if not called explicitly
    right angle length/.code={\def\rightanglelength{#1}},   % Length of symbol
    right angle length=2ex, % Make sure it is set...
    right angle symbol/.style n args={3}{
        insert path={
            let \p0 = ($(#1)!(#3)!(#2)$) in     % Intersection**
                let \p1 = ($(\p0)!\quadranta*\rightanglelength!(#3)$), % Point on base line
                \p2 = ($(\p0)!\quadrantb*\rightanglelength!(#2)$) in % Point on perpendicular line
                let \p3 = ($(\p1)+(\p2)-(\p0)$) in  % Corner point of symbol
            (\p1) -- (\p3) -- (\p2)
        }
    }
}
\pagestyle{fancy}
\fancyhf{}
\title{Boyle's Law Problem Set}
\author{Andi Jingzhi Tse, Janet Wan, Nancy Wang}
\date{January 16, 2023}
\lhead{Andi Jingzhi Tse, Janet Wan, Nancy Wang}
\rhead{January 16, 2023}
\pagestyle{fancy}

 
\cfoot{Page \thepage \hspace{1pt} of \pageref{LastPage}}


\begin{document}

\maketitle
\begin{custom-big}[Example 1]
List what experiment was involved in the founding of the Boyle’s law and explain how it works.
\end{custom-big}

\begin{custom-big}[Example 2]
Known that a balloon filled with 2 L of air will burst at a pressure of 3 atm. Which of the following condition will the balloon blow up. 
\begin{itemize}
    \item[(A)] 2 L of air with pressure of 2 atm
    \item[(B)] 3 L of air with pressure of 4 atm
    \item[(C)] 1 L of air with pressure of 4 atm
    \item[(D)] 4 L of air with pressure of 1 atm
\end{itemize}
\end{custom-big}

\begin{custom-big}[Example 3]
Write the equation for Boyle’s law and explain each of the variables.
\end{custom-big}


\begin{custom-big}[Example 4]
The formula of Boyle's law is $P_1V_2=P_2V_2$. What standard SI units are used in Boyle's Law?
\end{custom-big}

\begin{custom-big}[Example 5]
A gas occupies 10.6 liters at a pressure of 30.0 mmHg. What is the volume of the same gas when the pressure is increased to 70.0 mmHg?
\end{custom-big}

\begin{custom-big}[Example 6]

The pressure of a gas is 0.70 mmHg at a volume of 378.6 mL. When the pressure is lowered to 0.25 mmHg what is the volume in liters of the air?
\end{custom-big}



\begin{custom-big}[Example 7]
Janet arrived home and collapsed to her bed. The bed contained 10.0 L of air with pressure of 1.00 atm. The air is compressed to 3.00 L. Find the pressure of the air. 
\end{custom-big}

\begin{custom-big}[Example 8]
The gas initially has a volume of 5.0 L in the pump and a pressure of 1.0 atm. With the outlet of the pump plugged, if the pressure is changed to 3.5 atm by pressing down, what is the new volume?
\end{custom-big}


\begin{custom-big}[Example 9]
In a J-Tube the volume of air is measured 300 mL at a pressure of 24.5 mmHg. How much more mmHg should be added for the volume to lower to 250 mL?
\end{custom-big}


\begin{custom-big}[Example 10]
A sealed pot has volume of 4 L, 3.5 L of it is water, the rest is steam with pressure of 2.0 atm. 0.5 L of water is boiled into steam. Assume the volume of 1 atm steam is 10 times the volume of water with the same moles of water molecules. What is the pressure of steam after boiling?
\end{custom-big}


\begin{custom-big}[Challenge Question (Boyle's Law + Dalton's Law)]
Two bulbs of different volumes are separated by a valve. The valve between the 3.0 L bulb, in which the gas pressure is 2.0 atm, and the 2.0 L bulb, in which the gas pressure is 3.0 atm, is opened. What is the final pressure in the two bulbs, the temperature being constant and the same in both bulbs?
\end{custom-big}


\begin{custom-big}[Challenge Question (Boyle's Law + Dalton's Law)]
A glass will break if the pressure inside is more than 2.0 atm. The glass currently contains 3.7 L of air, with pressure of 1.3 atm. How many liters of air (1 atm) needs to be added to break the glass?
\end{custom-big}

\end{document}
